\subsection{注记与进一步结果}
\label{sec:ch9-notes-further-results}

\subsubsection{参考文献}
\label{subsec:ch9-references}

Cotlar 引理最早由 M. Cotlar 在自伴算子的情形中给出, 见 
\emph{A combinatorial inequality and its applications to $L^2$ spaces}, Rev. Mat. Cuyana 1 (1955), 41--55. 本章所述的一般情形由 M. Cotlar 与 E. M. Stein 各自独立得到, 最早见 A. Knapp 与 E. M. Stein 的论文 
\emph{Intertwining operators for semi-simple groups}, Ann. of Math. 93 (1971), 489--578. Carleson 测度及定理 \ref{thm:ch9-carleson-convolution-characterization} 中的刻画归功于 L. Carleson, 见 
\emph{An interpolation problem for bounded analytic functions}, Amer. J. Math. 80 (1958), 921--930, 以及 
\emph{Interpolation by bounded analytic functions and the corona problem}, Ann. of Math. 76 (1962), 547--559. Carleson 测度与 BMO 的关系归功于 C. Fefferman 与 E. M. Stein, 见 
\emph{$H^p$ spaces of several variables}, Acta Math. 129 (1972), 137--193.

$T1$ 定理归功于 G. David 与 J. L. Journé, 见 
\emph{A boundedness criterion for generalized Calderón--Zygmund operators}, Ann. of Math. 120 (1984), 371--397. 该定理有多种证明. 本章采用原始证明, 并参考 G. David 的学位论文中的进一步展开, Univ. Paris Sud, 1986. R. Coifman 与 Y. Meyer 给出了短得多的证明, 见 
\emph{A simple proof of a theorem by G. David and J.-L. Journé on singular integral operators}, 收于 \emph{Probability Theory and Harmonic Analysis}, J. Chao 与 W. Woyczyński 编, pp. 61--65, Marcel Dekker, New York, 1986. 这一证明也见于他们收入 Stein \MBUnresolvedReference{citation}{[16]} 的文章和 Torchinsky \MBUnresolvedReference{citation}{[19]} 的书. 原始证明以及另一种使用小波的证明见 Y. Meyer 的著作 \emph{Ondelettes et Opérateurs}, volume II, Hermann, Paris, 1990; 英译本见 Y. Meyer 与 R. Coifman, \emph{Wavelets: Calderón--Zygmund and Multilinear Operators}, Cambridge Univ. Press, Cambridge, 1997. 引理 \ref{lem:ch9-paraproduct-construction} 中的算子称为抛物积, 最早由 R. Coifman 与 Y. Meyer \MBUnresolvedReference{citation}{[2]} 引入. 关于其应用及其与 $T1$ 定理关系的综述, 见 M. Christ, \emph{Lectures on Singular Integral Operators}, CBMS Regional Conference Series in Mathematics 77, Amer. Math. Soc., Providence, 1990.

\hypertarget{chapter-09-page-225}{}

Stein \MBUnresolvedReference{citation}{[17]} 给出的 $T1$ 定理版本很有意思. 定理陈述中的弱有界性、$T1$ 和 $T^*1$ 条件被 $T$ 与 $T^*$ 的一个 ``限制有界性'' 条件代替:
\[
 \|T(\phi_{x,R})\|_2,\ \|T^*(\phi_{x,R})\|_2<CR^{n/2}.
\]
BMO 的作用出现在证明中: 限制有界性和核条件足以定义算子在常值函数上的作用, 并把它实现为一个 BMO 函数. 但 BMO 不出现在定理陈述中, 而陈述只涉及 $L^2$ 空间.

\subsubsection{Calderón 交换子与 Cauchy 积分}
\label{subsec:ch9-calderon-cauchy}

第 \ref{sec:ch9-t1-statement-applications} 节中作为 $T1$ 定理应用给出的算子最早由 A. P. Calderón 研究. 他证明了交换子 $T_1$ 在 $L^2$ 上有界, 见 
\emph{Commutators of singular integral operators}, Proc. Nat. Acad. Sci. U.S.A. 53 (1965), 1092--1099, 并证明了推论 \ref{cor:ch9-small-lipschitz-cauchy}, 见 
\emph{Cauchy integrals on Lipschitz curves and related operators}, Proc. Nat. Acad. Sci. U.S.A. 74 (1977), 1324--1327. 另见他的综述 
\emph{Commutators, singular integrals on Lipschitz curves and applications}, Proceedings of the I.C.M. (Helsinki, 1978), pp. 85--96, Acad. Sci. Fennica, Helsinki, 1978. R. Coifman 与 Y. Meyer 证明了其余交换子在 $L^2$ 上的有界性, 见 
\emph{Commutateurs d'intégrales singulières et opérateurs multilinéaires}, Ann. Inst. Fourier 28 (1978), 177--202, 但所得界比推论 \ref{cor:ch9-calderon-commutator-bounds} 中的界差, 其量级为 $k!\|A'\|_\infty^k$. 这些界不足以对推论 \ref{cor:ch9-small-lipschitz-cauchy} 证明中的级数求和.

推论 \ref{cor:ch9-small-lipschitz-cauchy} 中的限制 $\|A'\|_\infty<\varepsilon$ 并非必需; 对每个满足 $A'\in L^\infty$ 的 $A$, 结论都成立. R. Coifman、A. McIntosh 与 Y. Meyer 最先证明了这一点, 见 
\emph{L'intégrale de Cauchy définit un opérateur borné sur $L^2$ pour les courbes lipschitziennes}, Ann. of Math. 116 (1982), 361--387. G. David 从推论 \ref{cor:ch9-small-lipschitz-cauchy} 出发得到一般结果: 他证明, 若结论对 $\|A'\|_\infty\leq\varepsilon$ 成立, 则它对 $\|A'\|_\infty\leq\frac{10}{9}\varepsilon$ 成立, 见 
\emph{Opérateurs intégraux singuliers sur certaines courbes du plan complexe}, Ann. Sci. Éc. Norm. Sup. 17 (1984), 157--189. T. Murai 独立发展了类似论证, 见 
\emph{Boundedness of singular integral operators of Calderon type}, Proc. Japan Acad. Ser. A Math. Sci. 59 (1983), 364--367, 以及下文引用的著作. 此外, David 刻画了 Cauchy 积分在其上定义 $L^2$ 有界算子的曲线: 恰好是这样的曲线, 任意半径为 $r$ 的圆内部所含曲线片段的长度至多为 $Cr$, 其中 $C$ 固定. 这类曲线通常称为 Ahlfors--David 曲线.

T. Murai 的著作 \emph{A Real Variable Method for the Cauchy Transform and Analytic Capacity}, Lecture Notes in Math. 1307, Springer--Verlag, Berlin, 1988, 前两章研究交换子 $T_1$ 与 Cauchy 积分的 $L^2$ 有界性, 并给出这些结果的若干证明. 其中值得注意的是 P. Jones 与 S. Semmes 给出的 Cauchy 积分短证明.

\hypertarget{chapter-09-page-226}{}

该证明和另一个证明见这两位作者与 R. Coifman 合作的论文 
\emph{Two elementary proofs of the $L^2$ boundedness of Cauchy integrals on Lipschitz curves}, J. Amer. Math. Soc. 2 (1989), 553--564. 另见 P. Jones, 
\emph{Square functions, Cauchy integrals, analytic capacity, and harmonic measure}, 收于 \emph{Harmonic Analysis and Partial Differential Equations} (El Escorial, 1987), pp. 24--68, Lecture Notes in Math. 1384, Springer--Verlag, Berlin, 1989, 以及 S. Semmes, 
\emph{Square function estimates and the $T(b)$ theorem}, Proc. Amer. Math. Soc. 110 (1990), 721--726. M. Melnikov 与 J. Verdera 给出了一般结果的一个非常初等的几何证明, 见 
\emph{A geometric proof of the $L^2$ boundedness of the Cauchy integral on Lipschitz graphs}, Internat. Math. Res. Notices 1995, no. 7, 325--331.

\subsubsection{双线性 Hilbert 变换}
\label{subsec:ch9-bilinear-hilbert-transform}

在研究交换子 $T_1$ 的 $L^2$ 有界性时, Calderón 通过类似于旋转法的论证, 见第 \ref{sec:ch4-method-of-rotations} 节, 说明如果双线性奇异积分算子族
\[
 H_a(f,g)(x)=\lim_{\varepsilon\to0}
 \int_{|t|>\varepsilon}f(x-t)g(x+at)\,\frac{dt}{t}
\]
对 $-1\leq a\leq0$ 是从 $L^2(\mathbb R)\times L^\infty(\mathbb R)$ 到 $L^2(\mathbb R)$ 的一致有界映射, 那么 $T_1$ 有界. 当 $a=0$ 时, 该算子化为 $Hf\cdot g$; 当 $a=-1$ 时, 它化为 $H(f\cdot g)$, 其中 $H$ 是 Hilbert 变换. 因而问题在于证明 $H_a$ 对 $-1<a<0$ 一致有界.

算子 $H_1$ 称为双线性 Hilbert 变换. Calderón 还提出了一个较简单的问题: 证明 $H_1$ 是从 $L^2\times L^2$ 到 $L^1$ 的有界算子.

三十多年间, 这些问题没有取得进展. 随后 M. Lacey 与 C. Thiele 在两篇杰出的论文中证明了以下结果: 
\emph{$L^p$ estimates on the bilinear Hilbert transform for $2<p<\infty$}, Ann. of Math. 146 (1997), 693--724, 以及 
\emph{On Calderón's conjecture}, Ann. of Math. 149 (1999), 475--496.

\begin{Theorem}
\label{thm:ch9-bilinear-hilbert-transform}
给定 $a\in\mathbb R$ 以及指数 $p_1,p_2,p_3$, 满足 $1<p_1,p_2\leq\infty$, $p_3>2/3$ 且
\[
 \frac1{p_1}+\frac1{p_2}=\frac1{p_3},
\]
则
\[
 \|H_a(f,g)\|_{p_3}\leq C_{a,p_1,p_2}\|f\|_{p_1}\|g\|_{p_2}.
\]
\end{Theorem}

定理 \ref{thm:ch9-bilinear-hilbert-transform} 对双线性 Hilbert 变换确立了 Calderón 的猜想, 并表明算子 $H_a$, $-1<a<0$, 有界.

\hypertarget{chapter-09-page-227}{}

但是, 当 $a$ 趋于 $-1$ 或 $0$ 时, 常数 $C_{a,2,\infty}$ 无界, 因此该结果没有完成 Calderón 关于 $T_1$ 的 $L^2$ 有界性的证明. 本书付印时, L. Grafakos 与 X. Li 宣布已证明: 当 $2<p_1,p_2<\infty$ 且 $1<p_3<2$ 时, 该常数与 $a$ 无关; 对定理 \ref{thm:ch9-bilinear-hilbert-transform} 涵盖的所有其他情形, 只要 $a$ 与 $-1$ 保持正距离, 同样成立.

\subsubsection{$Tb$ 定理}
\label{subsec:ch9-tb-theorem}

$T1$ 定理不能给出比推论 \ref{cor:ch9-small-lipschitz-cauchy} 更强的 Cauchy 积分结论, 因为我们不能直接计算算子在函数 $1$ 上的作用. 然而, A. McIntosh 与 Y. Meyer 在 
\emph{Algèbres d'opérateurs définis par des intégrales singulières}, C. R. Acad. Sci. Paris 301 (1985), 395--397 中把 $T1$ 定理推广如下. 设函数 $b$ 有正下界; 更确切地说, 存在 $\delta>0$, 使得对所有 $y$, $\operatorname{Re}b(y)>\delta$. 若 $Tb=T^*b=0$, 且 $M_bTM_b$ 满足弱有界性, 其中 $M_b$ 表示乘以 $b$, 则 $T$ 在 $L^2$ 上有界. 令 $b(y)=1+iA'(y)$, 则 $b$ 有正下界, 并且可以证明 $b$ 的 Cauchy 积分为零, 因而该结果可直接用于 Cauchy 积分.

以这一结果为起点, G. David、J. L. Journé 与 S. Semmes 在 
\emph{Opérateurs de Calderón--Zygmund, fonctions para-accrétives et interpolation}, Rev. Mat. Iberoamericana 1 (4) (1985), 1--56 中得到 $T1$ 定理的一个重要推广, 即所谓的 $Tb$ 定理. 为陈述它, 需要几个定义. 给定 $\mu$, $0<\mu<1$, 令 $C_c^\mu(\mathbb R^n)$ 为满足
\[
 \|f\|_\mu=\sup_{x\neq y}\frac{|f(x)-f(y)|}{|x-y|^\mu}<\infty
\]
的紧支函数空间. 对函数 $b$, 令
\[
 bC_c^\mu=\{bf:f\in C_c^\mu\}.
\]
有界函数 $b:\mathbb R^n\to\mathbb C$ 称为拟增生的, 如果存在 $\delta>0$, 使得对每个 $x\in\mathbb R^n$ 和 $r>0$, 都存在球 $B=B(y,\rho)\subset B(x,r)$, 满足 $\rho\geq\delta r$ 且 $|b_B|\geq\delta$.

\begin{Theorem}
\label{thm:ch9-tb}
设 $b_1,b_2$ 是 $\mathbb R^n$ 上的两个拟增生函数, 且
\[
 T:b_1C_c^\mu(\mathbb R^n)\longrightarrow
 \bigl(b_2C_c^\mu(\mathbb R^n)\bigr)'
\]
是与标准核 $K(x,y)$ 相伴的算子. 则 $T$ 有连续的 $L^2$ 延拓, 当且仅当
\[
 Tb_1\in\operatorname{BMO},\qquad T^*b_2\in\operatorname{BMO},
\]
且 $M_{b_2}TM_{b_1}$ 满足弱有界性.
\end{Theorem}

\hypertarget{chapter-09-page-228}{}

由于定理 \ref{thm:ch9-tb} 用空间 $C_c^\mu$ 而非 $C_c^\infty$ 表述, 它可以推广到齐次型空间, 见第 \ref{subsec:ch5-spaces-homogeneous-type} 节. 该结果还有其他变体, 包括对不属于 $C^\infty$ 的函数 $b$ 定义 $Tb$, 以及采用稍有不同的弱有界性. 关于所有这些结果, 可参阅以上引用的论文. 另见 M. Christ, 
\emph{A $T(b)$ theorem with remarks on analytic capacity and the Cauchy integral}, Colloq. Math. 60/61 (1990), 601--628, 以及以上引用的 Meyer 著作.
